\subsection{Implementação do Banco de Dados}

A implementação do banco de dados do sistema Tropical Mix foi realizada com base em um modelo relacional, utilizando o MySQL como sistema gerenciador de banco de dados. A estrutura foi planejada para armazenar e relacionar as informações necessárias ao funcionamento da aplicação, como produtos, categorias, usuários, clientes, fornecedores, pedidos, itens vendidos, abastecimentos e movimentações de estoque.

O banco de dados atua como camada de persistência da aplicação, garantindo que as informações registradas pelo sistema sejam armazenadas de forma estruturada, consistente e acessível pelas demais camadas do projeto. Para isso, foram utilizadas chaves primárias, chaves estrangeiras e restrições de integridade, permitindo manter o relacionamento correto entre as tabelas e evitar registros inconsistentes.

Além do armazenamento dos dados cadastrais, a estrutura do banco também apoia regras importantes do sistema, como o controle da quantidade atual dos produtos, o registro de entradas e saídas de estoque, o vínculo entre pedidos e clientes, e a identificação do usuário responsável pelas operações realizadas.

\subsubsection{Dicionário de Dados}

O dicionário de dados apresenta a descrição das tabelas implementadas no banco de dados do sistema Tropical Mix. Para cada tabela, são descritos seus campos, tipos de dados, restrições e finalidades dentro da aplicação.

Essa documentação é importante porque facilita a compreensão da estrutura do banco de dados, apoia futuras manutenções e permite identificar como as informações do sistema estão organizadas. Dessa forma, o dicionário de dados funciona como uma referência técnica para compreender a relação entre as entidades utilizadas no controle de estoque, vendas e cadastros do sistema.

\subsubsection{Tabela: Categoria}

A tabela \textit{Categoria} é responsável pelo armazenamento das categorias dos produtos cadastrados no sistema. Sua finalidade é permitir a organização e classificação dos produtos comercializados pela empresa, facilitando processos de gerenciamento, consulta e controle de estoque.

\begin{table}[H]
\centering
\begin{tabular}{|p{3cm}|p{3cm}|p{3cm}|p{5cm}|}
\hline
\textbf{Campo} & \textbf{Tipo} & \textbf{Restrição} & \textbf{Descrição} \\ \hline
id\_categoria & INT & PK, AI & Identificador único da categoria \\ \hline
nome & VARCHAR(100) & NOT NULL & Nome da categoria \\ \hline
descricao & VARCHAR(255) & NULL & Descrição da categoria \\ \hline
\end{tabular}
\caption{Tabela: Categoria}
\end{table}

\subsubsection{Tabela: Produto}

A tabela \textit{Produto} armazena as informações relacionadas aos produtos comercializados pela empresa. Nela são registrados dados como nome, sabor, marca, quantidade atual, data de vencimento, preço unitário e categoria vinculada.

Essa tabela possui papel central no sistema, pois é utilizada tanto nas operações de estoque quanto nas vendas. A quantidade atual do produto permite acompanhar a disponibilidade dos itens, enquanto o vínculo com a tabela de categoria facilita a organização dos produtos e a aplicação de filtros na interface.

\begin{table}[H]
\centering
\begin{tabular}{|p{3cm}|p{3cm}|p{3cm}|p{5cm}|}
\hline
\textbf{Campo} & \textbf{Tipo} & \textbf{Restrição} & \textbf{Descrição} \\ \hline
id\_produto & INT & PK, AI & Identificador único do produto \\ \hline
nome\_produto & VARCHAR(150) & NOT NULL & Nome do produto \\ \hline
sabor & VARCHAR(50) & NULL & Sabor do produto \\ \hline
marca & VARCHAR(100) & NULL & Marca do produto \\ \hline
qtdd\_atual & INT & DEFAULT 0 & Quantidade disponível em estoque \\ \hline
data\_vencimento & DATE & NULL & Data de vencimento do produto \\ \hline
preco\_unitario & DECIMAL(10,2) & NOT NULL & Valor unitário do produto \\ \hline
dt\_cadastro & DATETIME & DEFAULT CURRENT\_TIMESTAMP & Data de cadastro do produto \\ \hline
id\_categoria & INT & FK, NOT NULL & Referência da categoria do produto \\ \hline
ativo & BOOLEAN & DEFAULT TRUE & Status de ativação do produto \\ \hline
\end{tabular}
\caption{Tabela: Produto}
\end{table}

\subsubsection{Tabela: Usuário}

A tabela \textit{Usuário} é responsável pelo armazenamento dos dados dos usuários do sistema. Essa estrutura permite o controle de autenticação, identificação e definição de níveis de acesso dos colaboradores que utilizam a aplicação.

Por meio dessa tabela, o sistema diferencia usuários administradores e vendedores, permitindo restringir funcionalidades conforme o perfil de acesso. Além disso, o usuário é relacionado a pedidos e movimentações de estoque, possibilitando identificar o responsável por determinadas operações.

\begin{table}[H]
\centering
\begin{tabular}{|p{3cm}|p{3cm}|p{3cm}|p{5cm}|}
\hline
\textbf{Campo} & \textbf{Tipo} & \textbf{Restrição} & \textbf{Descrição} \\ \hline
id\_usuario & INT & PK, AI & Identificador único do usuário \\ \hline
nome\_usuario & VARCHAR(100) & NOT NULL & Nome do usuário \\ \hline
cpf & VARCHAR(14) & NULL & CPF do usuário \\ \hline
email & VARCHAR(150) & UNIQUE, NOT NULL & E-mail do usuário \\ \hline
senha & VARCHAR(255) & NOT NULL & Senha de acesso armazenada em formato criptografado \\ \hline
nivel\_acesso & ENUM & DEFAULT 'vendedor' & Nível de acesso do usuário, podendo ser administrador ou vendedor \\ \hline
\end{tabular}
\caption{Tabela: Usuário}
\end{table}

\subsubsection{Tabela: Cliente}

A tabela \textit{Cliente} armazena as informações cadastrais dos clientes da empresa. Os dados registrados permitem identificar clientes vinculados às vendas, facilitando o controle comercial e o histórico de pedidos realizados.

Essa tabela é utilizada principalmente no processo de registro de pedidos, pois cada venda deve estar associada a um cliente. Dessa forma, o sistema consegue manter rastreabilidade das transações e consultar posteriormente quais clientes realizaram compras.

\begin{table}[H]
\centering
\begin{tabular}{|p{3cm}|p{3cm}|p{3cm}|p{5cm}|}
\hline
\textbf{Campo} & \textbf{Tipo} & \textbf{Restrição} & \textbf{Descrição} \\ \hline
id\_cliente & INT & PK, AI & Identificador único do cliente \\ \hline
razao\_social & VARCHAR(100) & NOT NULL & Nome ou razão social do cliente \\ \hline
cpf\_cnpj & VARCHAR(14) & NULL & CPF ou CNPJ do cliente \\ \hline
telefone & VARCHAR(20) & NULL & Telefone de contato \\ \hline
email & VARCHAR(150) & NULL & E-mail do cliente \\ \hline
endereco & VARCHAR(255) & NULL & Endereço do cliente \\ \hline
dt\_cadastro & DATETIME & DEFAULT CURRENT\_TIMESTAMP & Data de cadastro do cliente \\ \hline
\end{tabular}
\caption{Tabela: Cliente}
\end{table}

\subsubsection{Tabela: Pedido}

A tabela \textit{Pedido} representa o registro principal de uma venda realizada no sistema. Ela armazena informações gerais da transação, como o cliente relacionado, o usuário responsável, a data do pedido e o método de pagamento utilizado.

Essa tabela funciona como o cabeçalho da venda, pois reúne os dados principais do pedido, enquanto os produtos vendidos são detalhados na tabela \textit{Item Pedido}. O relacionamento com as tabelas \textit{Cliente} e \textit{Usuário} permite identificar tanto quem realizou a compra quanto o colaborador responsável pelo registro da operação.

\begin{table}[H]
\centering
\begin{tabular}{|p{3cm}|p{3cm}|p{3cm}|p{5cm}|}
\hline
\textbf{Campo} & \textbf{Tipo} & \textbf{Restrição} & \textbf{Descrição} \\ \hline
id\_pedido & INT & PK, AI & Identificador único do pedido \\ \hline
id\_cliente & INT & FK, NOT NULL & Referência ao cliente relacionado ao pedido \\ \hline
id\_usuario & INT & FK, NOT NULL & Referência ao usuário responsável pelo pedido \\ \hline
dt\_pedido & DATETIME & DEFAULT CURRENT\_TIMESTAMP & Data de realização do pedido \\ \hline
mtd\_pagamento & ENUM & NULL & Método de pagamento utilizado no pedido \\ \hline
\end{tabular}
\caption{Tabela: Pedido}
\end{table}

\subsubsection{Tabela: Tipo Movimentação}

A tabela \textit{Tipo Movimentação} armazena os tipos de movimentações possíveis no estoque, como entrada e saída. Sua finalidade é padronizar a classificação das operações realizadas sobre os produtos.

Essa tabela é utilizada em conjunto com a tabela \textit{Movimentação Estoque}, permitindo identificar se uma determinada movimentação corresponde a uma entrada de mercadoria, uma saída por venda ou outro tipo de alteração no estoque. Esse controle contribui para a rastreabilidade das operações e para a organização do histórico de movimentações.

\begin{table}[H]
\centering
\begin{tabular}{|p{3cm}|p{3cm}|p{3cm}|p{5cm}|}
\hline
\textbf{Campo} & \textbf{Tipo} & \textbf{Restrição} & \textbf{Descrição} \\ \hline
id\_tipo\_movimentacao & INT & PK, AI & Identificador único do tipo de movimentação \\ \hline
tipo\_movimentacao & VARCHAR(150) & NULL & Tipo de movimentação do estoque, como entrada ou saída \\ \hline
\end{tabular}
\caption{Tabela: Tipo Movimentação}
\end{table}

\subsubsection{Tabela: Movimentação Estoque}

A tabela \textit{Movimentação Estoque} registra o histórico das alterações realizadas na quantidade dos produtos. Cada registro representa uma movimentação de estoque, contendo o produto afetado, o usuário responsável, o tipo de movimentação, a quantidade movimentada e a data da atualização.

Essa tabela é essencial para a rastreabilidade do sistema, pois permite acompanhar entradas e saídas de produtos ao longo do tempo. Além de apoiar o controle operacional, o histórico de movimentações também pode ser utilizado para auditoria, conferência de estoque e geração de relatórios gerenciais.

\begin{table}[H]
\centering
\begin{tabular}{|p{3cm}|p{3cm}|p{3cm}|p{5cm}|}
\hline
\textbf{Campo} & \textbf{Tipo} & \textbf{Restrição} & \textbf{Descrição} \\ \hline
id\_estoque & INT & PK, AI & Identificador único da movimentação de estoque \\ \hline
id\_produto & INT & FK, NOT NULL & Referência ao produto movimentado \\ \hline
id\_usuario & INT & FK, NOT NULL & Referência ao usuário responsável pela movimentação \\ \hline
id\_tipo\_movimentacao & INT & FK, NOT NULL & Referência ao tipo de movimentação realizada \\ \hline
qtdd\_movimentacao & INT & DEFAULT 0 & Quantidade movimentada no estoque \\ \hline
ultima\_atualizacao & DATETIME & DEFAULT CURRENT\_TIMESTAMP & Data e hora da movimentação \\ \hline
\end{tabular}
\caption{Tabela: Movimentação Estoque}
\end{table}

\subsubsection{Tabela: Item Pedido}

A tabela \textit{Item Pedido} armazena os produtos associados a cada pedido realizado. Enquanto a tabela \textit{Pedido} registra os dados gerais da venda, a tabela \textit{Item Pedido} detalha quais produtos foram vendidos, suas quantidades, preços unitários e subtotal.

Essa estrutura permite que um mesmo pedido contenha vários produtos, mantendo a relação entre a venda e seus respectivos itens. O uso dessa tabela também facilita o cálculo do valor total da venda e a atualização do estoque com base nas quantidades vendidas.

\begin{table}[H]
\centering
\begin{tabular}{|p{3cm}|p{3cm}|p{3cm}|p{5cm}|}
\hline
\textbf{Campo} & \textbf{Tipo} & \textbf{Restrição} & \textbf{Descrição} \\ \hline
id\_item\_pedido & INT & PK, AI & Identificador único do item do pedido \\ \hline
id\_pedido & INT & FK, NOT NULL & Referência ao pedido relacionado \\ \hline
id\_produto & INT & FK, NOT NULL & Referência ao produto vendido \\ \hline
qtdd\_pedido & INT & NOT NULL & Quantidade solicitada do produto \\ \hline
preco\_unitario & DECIMAL(10,2) & NULL & Valor unitário do produto no momento da venda \\ \hline
subtotal & DECIMAL(10,2) & GENERATED & Valor total do item, calculado pela quantidade multiplicada pelo preço unitário \\ \hline
\end{tabular}
\caption{Tabela: Item Pedido}
\end{table}

\subsubsection{Tabela: Fornecedor}

A tabela \textit{Fornecedor} armazena os dados cadastrais dos fornecedores da empresa, como razão social, CNPJ, telefone e e-mail. Sua finalidade é permitir o controle das empresas ou pessoas responsáveis pelo fornecimento dos produtos comercializados pela Tropical Mix.

Essa tabela está diretamente relacionada ao processo de abastecimento do estoque, pois cada entrada de produto pode ser vinculada a um fornecedor. Dessa forma, o sistema mantém o histórico de origem dos produtos e facilita consultas relacionadas a compras, reposições e contatos comerciais.

\begin{table}[H]
\centering
\begin{tabular}{|p{3cm}|p{3cm}|p{3cm}|p{5cm}|}
\hline
\textbf{Campo} & \textbf{Tipo} & \textbf{Restrição} & \textbf{Descrição} \\ \hline
id\_fornecedor & INT & PK, AI & Identificador único do fornecedor \\ \hline
razao\_social & VARCHAR(150) & NOT NULL & Nome ou razão social do fornecedor \\ \hline
cnpj & VARCHAR(18) & UNIQUE & CNPJ do fornecedor \\ \hline
telefone & VARCHAR(20) & NULL & Telefone de contato do fornecedor \\ \hline
email & VARCHAR(150) & NULL & E-mail do fornecedor \\ \hline
\end{tabular}
\caption{Tabela: Fornecedor}
\end{table}

\subsubsection{Tabela: Abastecimento}

A tabela \textit{Abastecimento} registra as entradas de produtos no estoque. Cada registro indica o fornecedor responsável, o produto abastecido, a data da entrada, a quantidade recebida, a quantidade disponível no lote, o valor unitário e a data de vencimento do lote.

Essa tabela é utilizada para controlar as reposições de estoque e manter o histórico de abastecimentos realizados. Ao registrar uma nova entrada, o sistema consegue atualizar a quantidade disponível do produto e preservar informações importantes para acompanhamento de custos, fornecedores, lotes, validade e movimentações futuras.

\begin{table}[H]
\centering
\begin{tabular}{|p{3cm}|p{3cm}|p{3cm}|p{5cm}|}
\hline
\textbf{Campo} & \textbf{Tipo} & \textbf{Restrição} & \textbf{Descrição} \\ \hline
id\_abastecimento & INT & PK, AI & Identificador único do abastecimento \\ \hline
id\_fornecedor & INT & FK, NOT NULL & Referência ao fornecedor responsável pelo abastecimento \\ \hline
id\_produto & INT & FK, NOT NULL & Referência ao produto abastecido \\ \hline
numero\_lote & VARCHAR(50) & NULL & Número de identificação do lote do produto \\ \hline
dt\_abastecimento & DATETIME & DEFAULT CURRENT\_TIMESTAMP & Data e hora do abastecimento \\ \hline
qtdd\_recebida & INT & NOT NULL & Quantidade recebida no abastecimento \\ \hline
qtdd\_disponivel & INT & NULL & Quantidade ainda disponível no lote \\ \hline
valor\_unitario & DECIMAL(10,2) & NULL & Valor unitário do produto no abastecimento \\ \hline
data\_vencimento & DATE & NULL & Data de vencimento do lote \\ \hline
id\_usuario & INT & FK, NOT NULL & Usuário responsável pelo registro do abastecimento \\ \hline
\end{tabular}
\caption{Tabela: Abastecimento}
\end{table}

\subsubsection{Tabela: Categoria}

A tabela \textit{Categoria} é responsável pelo armazenamento das categorias dos produtos cadastrados no sistema. Sua finalidade é permitir a organização e classificação dos produtos comercializados pela empresa, facilitando processos de gerenciamento, consulta e controle de estoque.

\begin{table}[H]
\centering
\begin{tabular}{|p{3cm}|p{3cm}|p{3cm}|p{5cm}|}
\hline
\textbf{Campo} & \textbf{Tipo} & \textbf{Restrição} & \textbf{Descrição} \\ \hline
id\_categoria & INT & PK, AI & Identificador único da categoria \\ \hline
nome & VARCHAR(100) & NOT NULL & Nome da categoria \\ \hline
descricao & VARCHAR(255) & NULL & Descrição da categoria \\ \hline
\end{tabular}
\caption{Tabela: Categoria}
\end{table}

\subsubsection{Tabela: Produto}

A tabela \textit{Produto} armazena as informações relacionadas aos produtos comercializados pela empresa. Nela são registrados dados como nome, sabor, marca, quantidade atual, data de vencimento, preço unitário e categoria vinculada.

Essa tabela possui papel central no sistema, pois é utilizada tanto nas operações de estoque quanto nas vendas. A quantidade atual do produto permite acompanhar a disponibilidade dos itens, enquanto o vínculo com a tabela de categoria facilita a organização dos produtos e a aplicação de filtros na interface.

\begin{table}[H]
\centering
\begin{tabular}{|p{3cm}|p{3cm}|p{3cm}|p{5cm}|}
\hline
\textbf{Campo} & \textbf{Tipo} & \textbf{Restrição} & \textbf{Descrição} \\ \hline
id\_produto & INT & PK, AI & Identificador único do produto \\ \hline
nome\_produto & VARCHAR(150) & NOT NULL & Nome do produto \\ \hline
sabor & VARCHAR(50) & NULL & Sabor do produto \\ \hline
marca & VARCHAR(100) & NULL & Marca do produto \\ \hline
qtdd\_atual & INT & DEFAULT 0 & Quantidade disponível em estoque \\ \hline
data\_vencimento & DATE & NULL & Data de vencimento do produto \\ \hline
preco\_unitario & DECIMAL(10,2) & NOT NULL & Valor unitário do produto \\ \hline
dt\_cadastro & DATETIME & DEFAULT CURRENT\_TIMESTAMP & Data de cadastro do produto \\ \hline
id\_categoria & INT & FK, NOT NULL & Referência da categoria do produto \\ \hline
ativo & BOOLEAN & DEFAULT TRUE & Status de ativação do produto \\ \hline
\end{tabular}
\caption{Tabela: Produto}
\end{table}

\subsubsection{Tabela: Usuário}

A tabela \textit{Usuário} é responsável pelo armazenamento dos dados dos usuários do sistema. Essa estrutura permite o controle de autenticação, identificação e definição de níveis de acesso dos colaboradores que utilizam a aplicação.

Por meio dessa tabela, o sistema diferencia usuários administradores e vendedores, permitindo restringir funcionalidades conforme o perfil de acesso. Além disso, o usuário é relacionado a pedidos e movimentações de estoque, possibilitando identificar o responsável por determinadas operações.

\begin{table}[H]
\centering
\begin{tabular}{|p{3cm}|p{3cm}|p{3cm}|p{5cm}|}
\hline
\textbf{Campo} & \textbf{Tipo} & \textbf{Restrição} & \textbf{Descrição} \\ \hline
id\_usuario & INT & PK, AI & Identificador único do usuário \\ \hline
nome\_usuario & VARCHAR(100) & NOT NULL & Nome do usuário \\ \hline
cpf & VARCHAR(14) & NULL & CPF do usuário \\ \hline
email & VARCHAR(150) & UNIQUE, NOT NULL & E-mail do usuário \\ \hline
senha & VARCHAR(255) & NOT NULL & Senha de acesso armazenada em formato criptografado \\ \hline
nivel\_acesso & ENUM & DEFAULT 'vendedor' & Nível de acesso do usuário, podendo ser administrador ou vendedor \\ \hline
\end{tabular}
\caption{Tabela: Usuário}
\end{table}

\subsubsection{Tabela: Cliente}

A tabela \textit{Cliente} armazena as informações cadastrais dos clientes da empresa. Os dados registrados permitem identificar clientes vinculados às vendas, facilitando o controle comercial e o histórico de pedidos realizados.

Essa tabela é utilizada principalmente no processo de registro de pedidos, pois cada venda deve estar associada a um cliente. Dessa forma, o sistema consegue manter rastreabilidade das transações e consultar posteriormente quais clientes realizaram compras.

\begin{table}[H]
\centering
\begin{tabular}{|p{3cm}|p{3cm}|p{3cm}|p{5cm}|}
\hline
\textbf{Campo} & \textbf{Tipo} & \textbf{Restrição} & \textbf{Descrição} \\ \hline
id\_cliente & INT & PK, AI & Identificador único do cliente \\ \hline
razao\_social & VARCHAR(100) & NOT NULL & Nome ou razão social do cliente \\ \hline
cpf\_cnpj & VARCHAR(14) & NULL & CPF ou CNPJ do cliente \\ \hline
telefone & VARCHAR(20) & NULL & Telefone de contato \\ \hline
email & VARCHAR(150) & NULL & E-mail do cliente \\ \hline
endereco & VARCHAR(255) & NULL & Endereço do cliente \\ \hline
dt\_cadastro & DATETIME & DEFAULT CURRENT\_TIMESTAMP & Data de cadastro do cliente \\ \hline
\end{tabular}
\caption{Tabela: Cliente}
\end{table}

\subsubsection{Tabela: Pedido}

A tabela \textit{Pedido} representa o registro principal de uma venda realizada no sistema. Ela armazena informações gerais da transação, como o cliente relacionado, o usuário responsável, a data do pedido e o método de pagamento utilizado.

Essa tabela funciona como o cabeçalho da venda, pois reúne os dados principais do pedido, enquanto os produtos vendidos são detalhados na tabela \textit{Item Pedido}. O relacionamento com as tabelas \textit{Cliente} e \textit{Usuário} permite identificar tanto quem realizou a compra quanto o colaborador responsável pelo registro da operação.

\begin{table}[H]
\centering
\begin{tabular}{|p{3cm}|p{3cm}|p{3cm}|p{5cm}|}
\hline
\textbf{Campo} & \textbf{Tipo} & \textbf{Restrição} & \textbf{Descrição} \\ \hline
id\_pedido & INT & PK, AI & Identificador único do pedido \\ \hline
id\_cliente & INT & FK, NOT NULL & Referência ao cliente relacionado ao pedido \\ \hline
id\_usuario & INT & FK, NOT NULL & Referência ao usuário responsável pelo pedido \\ \hline
dt\_pedido & DATETIME & DEFAULT CURRENT\_TIMESTAMP & Data de realização do pedido \\ \hline
mtd\_pagamento & ENUM & NULL & Método de pagamento utilizado no pedido \\ \hline
\end{tabular}
\caption{Tabela: Pedido}
\end{table}

\subsubsection{Tabela: Tipo Movimentação}

A tabela \textit{Tipo Movimentação} armazena os tipos de movimentações possíveis no estoque, como entrada e saída. Sua finalidade é padronizar a classificação das operações realizadas sobre os produtos.

Essa tabela é utilizada em conjunto com a tabela \textit{Movimentação Estoque}, permitindo identificar se uma determinada movimentação corresponde a uma entrada de mercadoria, uma saída por venda ou outro tipo de alteração no estoque. Esse controle contribui para a rastreabilidade das operações e para a organização do histórico de movimentações.

\begin{table}[H]
\centering
\begin{tabular}{|p{3cm}|p{3cm}|p{3cm}|p{5cm}|}
\hline
\textbf{Campo} & \textbf{Tipo} & \textbf{Restrição} & \textbf{Descrição} \\ \hline
id\_tipo\_movimentacao & INT & PK, AI & Identificador único do tipo de movimentação \\ \hline
tipo\_movimentacao & VARCHAR(150) & NULL & Tipo de movimentação do estoque, como entrada ou saída \\ \hline
\end{tabular}
\caption{Tabela: Tipo Movimentação}
\end{table}

\subsubsection{Tabela: Movimentação Estoque}

A tabela \textit{Movimentação Estoque} registra o histórico das alterações realizadas na quantidade dos produtos. Cada registro representa uma movimentação de estoque, contendo o produto afetado, o usuário responsável, o tipo de movimentação, a quantidade movimentada e a data da atualização.

Essa tabela é essencial para a rastreabilidade do sistema, pois permite acompanhar entradas e saídas de produtos ao longo do tempo. Além de apoiar o controle operacional, o histórico de movimentações também pode ser utilizado para auditoria, conferência de estoque e geração de relatórios gerenciais.

\begin{table}[H]
\centering
\begin{tabular}{|p{3cm}|p{3cm}|p{3cm}|p{5cm}|}
\hline
\textbf{Campo} & \textbf{Tipo} & \textbf{Restrição} & \textbf{Descrição} \\ \hline
id\_estoque & INT & PK, AI & Identificador único da movimentação de estoque \\ \hline
id\_produto & INT & FK, NOT NULL & Referência ao produto movimentado \\ \hline
id\_usuario & INT & FK, NOT NULL & Referência ao usuário responsável pela movimentação \\ \hline
id\_tipo\_movimentacao & INT & FK, NOT NULL & Referência ao tipo de movimentação realizada \\ \hline
qtdd\_movimentacao & INT & DEFAULT 0 & Quantidade movimentada no estoque \\ \hline
ultima\_atualizacao & DATETIME & DEFAULT CURRENT\_TIMESTAMP & Data e hora da movimentação \\ \hline
\end{tabular}
\caption{Tabela: Movimentação Estoque}
\end{table}

\subsubsection{Tabela: Item Pedido}

A tabela \textit{Item Pedido} armazena os produtos associados a cada pedido realizado. Enquanto a tabela \textit{Pedido} registra os dados gerais da venda, a tabela \textit{Item Pedido} detalha quais produtos foram vendidos, suas quantidades, preços unitários e subtotal.

Essa estrutura permite que um mesmo pedido contenha vários produtos, mantendo a relação entre a venda e seus respectivos itens. O uso dessa tabela também facilita o cálculo do valor total da venda e a atualização do estoque com base nas quantidades vendidas.

\begin{table}[H]
\centering
\begin{tabular}{|p{3cm}|p{3cm}|p{3cm}|p{5cm}|}
\hline
\textbf{Campo} & \textbf{Tipo} & \textbf{Restrição} & \textbf{Descrição} \\ \hline
id\_item\_pedido & INT & PK, AI & Identificador único do item do pedido \\ \hline
id\_pedido & INT & FK, NOT NULL & Referência ao pedido relacionado \\ \hline
id\_produto & INT & FK, NOT NULL & Referência ao produto vendido \\ \hline
qtdd\_pedido & INT & NOT NULL & Quantidade solicitada do produto \\ \hline
preco\_unitario & DECIMAL(10,2) & NULL & Valor unitário do produto no momento da venda \\ \hline
subtotal & DECIMAL(10,2) & GENERATED & Valor total do item, calculado pela quantidade multiplicada pelo preço unitário \\ \hline
\end{tabular}
\caption{Tabela: Item Pedido}
\end{table}

\subsubsection{Tabela: Fornecedor}

A tabela \textit{Fornecedor} armazena os dados cadastrais dos fornecedores da empresa, como razão social, CNPJ, telefone e e-mail. Sua finalidade é permitir o controle das empresas ou pessoas responsáveis pelo fornecimento dos produtos comercializados pela Tropical Mix.

Essa tabela está diretamente relacionada ao processo de abastecimento do estoque, pois cada entrada de produto pode ser vinculada a um fornecedor. Dessa forma, o sistema mantém o histórico de origem dos produtos e facilita consultas relacionadas a compras, reposições e contatos comerciais.

\begin{table}[H]
\centering
\begin{tabular}{|p{3cm}|p{3cm}|p{3cm}|p{5cm}|}
\hline
\textbf{Campo} & \textbf{Tipo} & \textbf{Restrição} & \textbf{Descrição} \\ \hline
id\_fornecedor & INT & PK, AI & Identificador único do fornecedor \\ \hline
razao\_social & VARCHAR(150) & NOT NULL & Nome ou razão social do fornecedor \\ \hline
cnpj & VARCHAR(18) & UNIQUE & CNPJ do fornecedor \\ \hline
telefone & VARCHAR(20) & NULL & Telefone de contato do fornecedor \\ \hline
email & VARCHAR(150) & NULL & E-mail do fornecedor \\ \hline
\end{tabular}
\caption{Tabela: Fornecedor}
\end{table}

\subsubsection{Tabela: Abastecimento}

A tabela \textit{Abastecimento} registra as entradas de produtos no estoque. Cada registro indica o fornecedor responsável, o produto abastecido, a data da entrada, a quantidade recebida, a quantidade disponível no lote, o valor unitário e a data de vencimento do lote.

Essa tabela é utilizada para controlar as reposições de estoque e manter o histórico de abastecimentos realizados. Ao registrar uma nova entrada, o sistema consegue atualizar a quantidade disponível do produto e preservar informações importantes para acompanhamento de custos, fornecedores, lotes, validade e movimentações futuras.

\begin{table}[H]
\centering
\begin{tabular}{|p{3cm}|p{3cm}|p{3cm}|p{5cm}|}
\hline
\textbf{Campo} & \textbf{Tipo} & \textbf{Restrição} & \textbf{Descrição} \\ \hline
id\_abastecimento & INT & PK, AI & Identificador único do abastecimento \\ \hline
id\_fornecedor & INT & FK, NOT NULL & Referência ao fornecedor responsável pelo abastecimento \\ \hline
id\_produto & INT & FK, NOT NULL & Referência ao produto abastecido \\ \hline
numero\_lote & VARCHAR(50) & NULL & Número de identificação do lote do produto \\ \hline
dt\_abastecimento & DATETIME & DEFAULT CURRENT\_TIMESTAMP & Data e hora do abastecimento \\ \hline
qtdd\_recebida & INT & NOT NULL & Quantidade recebida no abastecimento \\ \hline
qtdd\_disponivel & INT & NULL & Quantidade ainda disponível no lote \\ \hline
valor\_unitario & DECIMAL(10,2) & NULL & Valor unitário do produto no abastecimento \\ \hline
data\_vencimento & DATE & NULL & Data de vencimento do lote \\ \hline
id\_usuario & INT & FK, NOT NULL & Usuário responsável pelo registro do abastecimento \\ \hline
\end{tabular}
\caption{Tabela: Abastecimento}
\end{table}